% Traditional Chinese PDF settings for Romans Deep Study.
\usepackage{fontspec}
\usepackage{xeCJK}
\setmainfont{Times New Roman}
\newfontfamily\greekfallback{Arial Unicode MS}
% STSong is available in the current TeX/macOS publishing environment. A
% Linux CI runner should provision a licensed CJK font and verify the embedded
% font list in the PDF before release.
\setCJKmainfont{STSong}
\setCJKsansfont{STSong}
\setCJKmonofont{STSong}
\usepackage{xcolor}
\definecolor{ScriptureGold}{HTML}{8A641C}
\usepackage{microtype}
\usepackage{geometry}
\geometry{letterpaper,margin=24mm}
\setlength{\parindent}{0pt}
\setlength{\parskip}{0.35em}
\linespread{1.08}
\clubpenalty=10000
\widowpenalty=10000
\emergencystretch=2em
\usepackage{fancyhdr}
\setlength{\headheight}{16pt}
\pagestyle{fancy}
\fancyhf{}
\fancyhead[L]{\small\nouppercase{\leftmark}}
\fancyhead[R]{\small\thepage}
\renewcommand{\headrulewidth}{0pt}
\renewcommand{\chaptermark}[1]{\markboth{羅馬書研讀 · Romans Deep Study}{}}
\AtBeginDocument{\hypersetup{colorlinks=true,linkcolor=ScriptureGold,urlcolor=ScriptureGold,bookmarksopen=true}}

% A restrained, vector-only title page keeps the canonical PDF portable while
% making the assigned edition unmistakable in print and in thumbnails.
\makeatletter
\renewcommand{\maketitle}{%
  \begin{titlepage}
    \thispagestyle{empty}
    \begin{center}
      \vspace*{18mm}
      {\small\textcolor{ScriptureGold}{Three Books Deep Reading · Romans}\par}
      \vspace{2mm}
      {\small 三書精讀 · 羅馬書深度研讀\par}
      \vspace{18mm}
      {\color{ScriptureGold}\rule{0.48\textwidth}{0.6pt}\par}
      \vspace{10mm}
      {\fontsize{34}{40}\selectfont\bfseries\textcolor{ScriptureGold}{神的義}\par}
      {\Large\itshape The Righteousness of God\par}
      \vspace{6mm}
      {\fontsize{25}{31}\selectfont\bfseries 羅馬書研讀\par}
      {\large\itshape Romans Deep Study\par}
      \vspace{2mm}
      {\small\textcolor{ScriptureGold}{2026 整編版 · 2026 Edition}\par}
      \vspace{10mm}
      \fbox{%
        \begin{minipage}{0.78\textwidth}
          \centering
          \textcolor{ScriptureGold}{\large 「因為神的義正在這福音上顯明出來；這義是本於信，以致於信。」}\par
          \vspace{3mm}
          {\small\itshape “For in it the righteousness of God is revealed from faith to faith.”}\par
          \vspace{2mm}
          {\small 羅馬書 1:17 · Romans 1:17}
        \end{minipage}}\par
      \vfill
      {\small PubHub 三書精讀系統 · Soli Deo Gloria\par}
    \end{center}
  \end{titlepage}
}
\makeatother
